\chapter{标签：你预测什么，决定你能得到什么}\label{ch:label}

\begin{keybox}[本章要回答的问题]
为什么"预测明天涨跌"是一个坏问题？横截面排序、波动率归一化、三重障碍法各自
适合什么场景？重叠标签为什么会让你的 t 值虚高？A 股的 T+1 与涨跌停如何改变标签设计？
\end{keybox}

\section{一个被低估的事实：多数模型失败在标签，而不是模型}

研究新手最常见的路径是：拿到数据 $\rightarrow$ 构造几十个因子 $\rightarrow$ 上 LightGBM $\rightarrow$ 调参 $\rightarrow$
发现样本外不灵 $\rightarrow$ 换更复杂的模型。

他们很少检查一件事：\textbf{标签本身是否可预测}。

日频个股收益的信噪比极低。以本书的数据为例，某只股票某日的收益，
其横截面标准差约为 $2.5\%$；而任何模型能解释的部分，通常不到这个数字的 $5\%$。
在这种信噪比下，标签的设计方式（预测什么、怎么归一化、如何处理重叠）
对最终结果的影响，往往大于模型选择。

\section{三种标签形态}

\subsection{回归：预测未来 $h$ 日收益}

\begin{equation}
  y_{i,t} = \frac{\tilde P_{i,t+1+h}}{\tilde P_{i,t+1}} - 1
  \label{eq:label-ret}
\end{equation}

其中 $\tilde P$ 是后复权价，$t+1$ 表示"次日开盘买入"。

\textbf{优点}：保留了收益的幅度信息，与最终的组合权重直接对应。
\textbf{缺点}：厚尾。A 股单日涨跌 $20\%$ 的样本不少，少数极端值会主导平方损失。

\begin{warnbox}[一个常见错误：用收盘价构造标签]
\texttt{fwd\_ret = close[t+h]/close[t] - 1} 隐含了"$t$ 日收盘买入"的假设。
但你的信号是用 $t$ 日收盘数据算出来的，实际上根本来不及在 $t$ 日收盘成交。
正确做法是 \texttt{close[t+h]/open[t+1] - 1}——\textbf{信号与成交之间必须隔一个可执行的时点}。
\end{warnbox}

\subsection{横截面排序：预测相对强弱}

\begin{equation}
  y_{i,t} = \mathrm{rank}_{i}\!\left(r_{i,t+1\rightarrow t+1+h}\right) - \tfrac12
  \;\in\; \left(-\tfrac12,\ \tfrac12\right]
  \label{eq:label-rank}
\end{equation}

\textbf{优点}：
\begin{itemize}
  \item 天然抗极值（排名不受 $+50\%$ 与 $+5\%$ 的差距影响）；
  \item 与选股目标一致——我们只关心"买哪些"，不关心"市场整体涨多少"；
  \item 自动剔除市场 beta，模型学到的纯粹是横截面信息。
\end{itemize}

\textbf{缺点}：丢失幅度信息，无法区分"小幅跑赢"与"大幅跑赢"。

本书的主标签采用这一形式。这是业界的标准选择，原因是：日频收益的幅度部分
几乎全是噪声，而排序部分是相对稳定的。

\subsection{分类：涨 / 跌 / 前 $30\%$}

把问题变成二分类或多分类。

\textbf{优点}：可以直接用分类模型，输出概率有明确含义。
\textbf{缺点}：阈值是人为的；信息损失最大；概率校准需要额外工作。

除非你有明确的"只买前 $30\%$"的业务约束，否则不建议作为主标签，
可以用来做辅助诊断（例如分析模型在极端分位上的表现）。

\section{波动率归一化：让不同股票可比}

同样是 $+3\%$ 的周收益，对一只年化波动 $20\%$ 的银行股是"大涨"，
对一只年化波动 $80\%$ 的题材股是"平常"。

\begin{equation}
  y_{i,t} = \frac{r_{i,t+1\rightarrow t+1+h}}{\sigma_{i,t}\sqrt{h}},
  \qquad \sigma_{i,t} = \mathrm{std}(r_{i,t-59\rightarrow t})
  \label{eq:label-vol}
\end{equation}

除以 $\sigma$ 之后，标签近似服从同分布，模型的残差也更接近同方差。
\textbf{代价}：低波动股票的权重被放大，组合会向低波倾斜；
如果你同时用了低波因子，就会产生双重暴露。

\section{三重障碍法：什么时候值得用}

López de Prado 提出的三重障碍法（Triple Barrier）给每个样本定义三个出口：

\begin{itemize}
  \item 上障碍：$P_{i,t}(1 + u\sigma_i)$，止盈；
  \item 下障碍：$P_{i,t}(1 - d\sigma_i)$，止损；
  \item 垂直障碍：$t + h$ 个交易日，到期平仓。
\end{itemize}

标签取\textbf{先被触及的那个障碍}：上障碍 $+1$，下障碍 $-1$，到期 $0$。

\subsection{它解决了什么}

它把"路径依赖"纳入了标签：一个先跌 $8\%$ 再涨 $5\%$ 的样本，
在固定持有期标签下是 $+5\%$，在三重障碍下是 $-1$（先触发止损）。
如果你的实际交易规则带止损，三重障碍标签与真实盈亏更一致。

\subsection{它的代价}

\begin{enumerate}
  \item \textbf{标签不平衡}：多数样本落在 $0$（到期平仓），分类器会退化为常数预测。
  \item \textbf{参数敏感}：$u,d$ 的选择直接影响标签分布，而你又会用它来调模型——
        这是双重过拟合的入口。
  \item \textbf{重叠更严重}：路径依赖让相邻样本的标签相关性更高，有效样本数进一步下降。
  \item \textbf{计算量大}：需要对每个样本扫描未来 $h$ 天的路径。
\end{enumerate}

\begin{keybox}[本书的判断]
\textbf{日频选股不需要三重障碍法。}
它的设计目标是"交易执行"层面的路径依赖，更适合高频或单标的择时。
日频横截面选股的核心是"排序准不准"，固定持有期 + 横截面排序标签已经足够，
而且统计性质更干净。
\end{keybox}

\section{Meta-labeling：把方向和仓位分开}

Meta-labeling 的思路是两级模型：

\begin{enumerate}
  \item \textbf{主模型}决定方向（买 / 不买），可以是一条简单规则或一个粗模型；
  \item \textbf{副模型}预测"主模型这次判断对不对"，输出一个概率 $p$；
  \item 最终仓位 $\propto (2p - 1)$ 或只保留 $p > \theta$ 的样本。
\end{enumerate}

它的价值在于\textbf{降低假阳性}：主模型给出 100 个买入信号，副模型只保留其中 40 个
把握较大的。这在信号本身很稀疏（比如事件驱动）时特别有效。

在日频全市场选股中，因为每天都有几千个候选，"稀疏"不是问题，
meta-labeling 的边际收益有限，反而增加了复杂度。

\section{重叠标签：为什么你的 $t$ 值虚高}

这是本书最重要的一节之一。

假设你用 $h=5$ 的持有期标签，并且\textbf{每个交易日}都训练一个样本。那么：
\begin{itemize}
  \item $t$ 日的标签覆盖 $[t+1, t+6]$；
  \item $t+1$ 日的标签覆盖 $[t+2, t+7]$；
  \item 两者有 4 天的重叠。
\end{itemize}

后果是相邻样本高度相关。若你用普通 $t$ 检验判断 IC 是否显著，
标准误会严重低估，$t$ 值被夸大 $2$--$4$ 倍。

\begin{proposition}[有效样本数]
若标签持有期为 $h$ 个交易日、采样间隔为 $1$ 天，则\textbf{独立}样本数约为
观测数的 $1/h$，即
\begin{equation}
  n_{\text{eff}} \approx \frac{n}{h}.
\end{equation}
更精细的估计可以用 IC 序列的自相关来修正：
\begin{equation}
  n_{\text{eff}} = \frac{n}{1 + 2\sum_{k=1}^{h-1}\left(1-\frac{k}{h}\right)\rho_k},
\end{equation}
其中 $\rho_k$ 是 IC 序列的 $k$ 阶自相关。
\end{proposition}

工程上有两个正解：

\begin{enumerate}
  \item \textbf{间隔采样}：每 $h$ 个交易日取一个样本，让标签不重叠。
        本书的模型训练采用这一做法（每 5 个交易日一个样本）。
  \item \textbf{Newey-West 修正}：保留全部样本，但在算 $t$ 值时用 HAC 标准误，
        滞后阶数取 $\ge h-1$。本书的因子评价采用这一做法。
\end{enumerate}

\section{A 股的特殊约束}

\subsection{T+1}

当日买入不能当日卖出，因此：
\begin{itemize}
  \item 最短持有期是 1 个交易日（今天买、明天卖）；
  \item 标签的起点必须是 $t+1$（次日），不能是 $t$。
\end{itemize}

\subsection{涨跌停}

如果你的标签是"次日开盘到第 $h$ 日开盘的收益"，那么：
\begin{itemize}
  \item 次日一字涨停买不进 $\rightarrow$ 该样本应当被标记为不可执行；
  \item 到期日一字跌停卖不出 $\rightarrow$ 该样本的收益实际上无法实现。
\end{itemize}

这两种情况都集中在极端样本上，而极端样本恰恰是反转类因子最关注的对象。
不处理它们，反转因子的回测收益会被显著高估。

本书的处理方式：在样本层面保留，但在\textbf{组合构建}时剔除
\texttt{can\_buy\_next = False} 的股票（见第~\ref{ch:portfolio} 章）。

\section{标签与持有期的匹配}

持有期 $h$ 不是一个可以随便设的参数，它由三个因素共同决定：

\begin{enumerate}
  \item \textbf{信号的衰减速度}。反转信号 5 日内衰减大半，20 日基本消失；
        质量因子可以持有 3--6 个月。
  \item \textbf{交易成本}。$h$ 越小，换手越高，成本吃掉的比例越大。
        在单边 $30$ bp 的成本下，$h=1$ 的策略需要极高的 IC 才能覆盖成本。
  \item \textbf{容量}。$h$ 越小，策略容量越小。
\end{enumerate}

\begin{center}
\begin{tabular}{@{}lccc@{}}
\toprule
持有期 $h$ & 年换手（倍） & 单边成本要求 & 适合的信号类型 \\
\midrule
1 日 & 250 & 极高 & 微观结构（个人基本不可行） \\
5 日 & 50 & 高 & 短期反转、量价 \\
20 日 & 12 & 中 & 动量、波动、质量 \\
60 日 & 4 & 低 & 价值、质量、成长 \\
\bottomrule
\end{tabular}
\end{center}

本书选择 $h=5$ 作为主标签：它在"信号强度"与"成本可承受性"之间取得平衡，
并且能把日频数据的重叠问题控制在可管理范围内。

\section{本章小结}

\begin{enumerate}
  \item 标签设计比模型选择更重要，日频收益的信噪比极低。
  \item 横截面排序标签是选股任务的标准选择；回归标签保留幅度但受厚尾干扰。
  \item 波动率归一化让不同股票可比，但会引入低波暴露。
  \item 三重障碍法解决路径依赖，代价是标签不平衡与参数敏感；日频选股不必使用。
  \item 重叠标签会让 $t$ 值虚高，必须用间隔采样或 Newey-West 修正。
  \item A 股的 T+1 与涨跌停改变了标签的可执行性，必须在样本与组合两层处理。
\end{enumerate}
